\documentclass[tikz,border=8pt]{standalone}
\usepackage{tikz}
\usepackage{amsmath}
\usepackage{amssymb}
\usepackage{pgfplots}
\pgfplotsset{compat=1.18}
\usepackage{ctex}
\usetikzlibrary{calc,positioning,arrows.meta,matrix}

% LC 49 Group Anagrams：排序键哈希分组
% 输入 ["eat","tea","tan","ate","nat","bat"]

\begin{document}
\begin{tikzpicture}[
  every node/.style={font=\small},
  wordbox/.style={draw=gray!70, rectangle, rounded corners=2pt,
    minimum width=1.2cm, minimum height=0.6cm, thick, fill=white},
  keybox/.style={draw=blue!50, rectangle, rounded corners=2pt,
    minimum width=1.4cm, minimum height=0.6cm, thick, fill=blue!8},
  groupA/.style={draw=orange!60, rectangle, rounded corners=3pt,
    minimum width=1.2cm, minimum height=0.6cm, thick, fill=orange!8},
  groupB/.style={draw=purple!60, rectangle, rounded corners=3pt,
    minimum width=1.2cm, minimum height=0.6cm, thick, fill=purple!8},
  groupC/.style={draw=teal!60, rectangle, rounded corners=3pt,
    minimum width=1.2cm, minimum height=0.6cm, thick, fill=teal!8},
  htrow/.style={draw=gray!40, rectangle, rounded corners=2pt,
    minimum width=5.0cm, minimum height=0.65cm, thin, fill=gray!5},
]

%% ── 标题 ──────────────────────────────────────────────────
\node[font=\large\bfseries] at (7.5, 1.2)
  {字母异位词分组（LC 49 Group Anagrams）};
\node[font=\small,gray] at (7.5, 0.55)
  {key = 排序后字符串；value = 原词列表（哈希分桶）};

%% ── 输入数组 ──────────────────────────────────────────────
\node[font=\scriptsize\bfseries,gray] at (1.4, -0.4) {输入：};
\foreach \i/\w in {0/eat, 1/tea, 2/tan, 3/ate, 4/nat, 5/bat} {
  \node[wordbox,font=\ttfamily\small] at (2.5+\i*1.5, -0.4) {\w};
}

%% ── 排序演示（中间步骤）──────────────────────────────────
\node[font=\scriptsize\bfseries,gray] at (1.4, -1.5) {排序键：};
\foreach \i/\w/\k in {0/eat/aet, 1/tea/aet, 2/tan/ant, 3/ate/aet, 4/nat/ant, 5/bat/abt} {
  \node[keybox,font=\ttfamily\scriptsize] at (2.5+\i*1.5, -1.5) {\k};
  % 下箭头
  \draw[->,>=Stealth,gray!50,thin] (2.5+\i*1.5, -0.72) -- (2.5+\i*1.5, -1.22);
}

%% ── 哈希表 ────────────────────────────────────────────────
\node[font=\scriptsize\bfseries,gray,anchor=east] at (2.2, -3.0) {哈希表：};

% key "aet"
\node[keybox,font=\ttfamily] (ka) at (3.5, -2.8) {aet};
\node[font=\scriptsize] at (4.5, -2.8) {$\to$};
\node[groupA,font=\ttfamily] (wa1) at (5.6, -2.8) {eat};
\node[groupA,font=\ttfamily] (wa2) at (7.0, -2.8) {tea};
\node[groupA,font=\ttfamily] (wa3) at (8.4, -2.8) {ate};

% key "ant"
\node[keybox,font=\ttfamily] (kb) at (3.5, -3.7) {ant};
\node[font=\scriptsize] at (4.5, -3.7) {$\to$};
\node[groupB,font=\ttfamily] (wb1) at (5.6, -3.7) {tan};
\node[groupB,font=\ttfamily] (wb2) at (7.0, -3.7) {nat};

% key "abt"
\node[keybox,font=\ttfamily] (kc) at (3.5, -4.6) {abt};
\node[font=\scriptsize] at (4.5, -4.6) {$\to$};
\node[groupC,font=\ttfamily] (wc1) at (5.6, -4.6) {bat};

% 哈希表边框
\draw[draw=gray!40,rounded corners=4pt,dashed]
  (2.8,-2.4) rectangle (9.5,-5.1);
\node[font=\tiny,gray,anchor=north east] at (9.5,-2.4) {HashMap};

%% ── 引线：输入词 → 排序键行 ──────────────────────────────
% eat→aet行 / tea→aet行 / ate→aet行（均指向 ka 附近，避免穿越）
\draw[->,>=Stealth,orange!50,very thin,dashed]
  (2.5, -0.72) to[bend right=15] (3.5, -2.55);
\draw[->,>=Stealth,orange!50,very thin,dashed]
  (4.0, -0.72) to[bend right=10] (3.5, -2.55);
\draw[->,>=Stealth,orange!50,very thin,dashed]
  (7.0, -0.72) to[bend left=10] (3.5, -2.55);

\draw[->,>=Stealth,purple!50,very thin,dashed]
  (5.5, -0.72) to[bend right=5] (3.5, -3.45);
\draw[->,>=Stealth,purple!50,very thin,dashed]
  (8.5, -0.72) to[bend left=10] (3.5, -3.45);

\draw[->,>=Stealth,teal!50,very thin,dashed]
  (10.0, -0.72) to[bend left=15] (3.5, -4.35);

%% ── 最终分组结果 ──────────────────────────────────────────
\node[font=\scriptsize\bfseries,gray] at (1.4, -6.1) {输出：};

% 分组 A
\node[draw=orange!60,fill=orange!8,rounded corners=3pt,
      font=\small,inner sep=6pt,align=center]
  at (3.8, -6.1) {[eat, tea, ate]};

% 分组 B
\node[draw=purple!60,fill=purple!8,rounded corners=3pt,
      font=\small,inner sep=6pt,align=center]
  at (7.4, -6.1) {[tan, nat]};

% 分组 C
\node[draw=teal!60,fill=teal!8,rounded corners=3pt,
      font=\small,inner sep=6pt,align=center]
  at (10.6, -6.1) {[bat]};

%% ── 分隔线 ──────────────────────────────────────────────
\draw[gray!25,dashed] (0.0, -7.0) -- (15.0, -7.0);

%% ── 算法说明 ─────────────────────────────────────────────
\node[draw=gray!50,fill=white,rounded corners=3pt,font=\small,
      inner sep=8pt,align=left]
  at (7.5, -8.0)
  {\textbf{Group Anagrams 算法（时间 $O(nk\log k)$，空间 $O(nk)$）：}\\[4pt]
   对每个字符串 \texttt{s}：将其字符排序得到键 \texttt{key = sorted(s)}\\[2pt]
   以 \texttt{key} 为哈希表键，将 \texttt{s} 追加到对应值列表\\[2pt]
   遍历完成后，哈希表各 value 即为一个异位词组};

\end{tikzpicture}
\end{document}
